% [AUTHOR REWRITE] magnitudes
In Partial Pooling and Cost-Deterred Hiding the threshold equals the rating premium $b$. Because $\hat\beta_1$ measures the quality premium in log points against a local price per m\textsuperscript{2}, the nominal threshold $\bar R_m=\hat b_m\times P_m$ scales with the local price level, and high-price urban markets have the largest verified-quality premia in euros even when the proportional premium is similar across locations. Across the \attNMarketsApt{} estimable apartment markets, $\bar R_m$ averages \euro{}\attScBaseRbarApt{}/m\textsuperscript{2} (median \euro{}\attBaseRbarMedApt{}, range \euro{}\attBaseRbarMinApt{} to \euro{}\attBaseRbarMaxApt{}); across the \attNMarketsHouse{} house markets the mean is \euro{}\attScBaseRbarHouse{} and the median \euro{}\attBaseRbarMedHouse{} (Figure~\ref{fig:policy_rbar_dist}).\footnote{These figures use the rating premium itself, $\bar R=\hat b$. The exact-pricing solver of Section~\ref{sec:equilibrium}, which takes the threshold from the full-attention equilibrium, gives means of \euro{}\attRbarCoreEurTwoFiftyApt{} and \euro{}\attRbarCoreEurTwoFiftyHouse{} per m\textsuperscript{2}. The difference is consistent with linear versus exact pricing and not with attention, since it is of the same size in the two segments (Appendix~\ref{app:attention}).} Cross-market dispersion is substantial, so the markets at the upper end are intrinsically more amenable to privately funded renovation than those at the lower end.
